\exnumber{04}
\extitle{UX IV - Ex 5 - Prototyping}
\exfiles{ex04/Figma file, named Name\_FirstName\_DeliverableName\_Date\_VersionNumber}
\exauthorize{Those named in the exercise}
\makeheaderfiles

\textbf{Instructions:}

Give another dimension to your user flow and turn your user flow into a paper prototype (12 screens max).

First, draw low-fidelity prototypes on paper. One main screen = 1 mobile frame.

Then, draw your mid-fidelity screens on Figma with simple shapes, in black and white.

\textbf{Tips:} Vocabulary

\begin{itemize}
\item Low-Fidelity: Paper Prototypes - Example here \url{https://miro.medium.com/max/1400/1*5nUfqqA2gjdAYHagjbHA5w.jpeg}
\item Mid-Fidelity: Wireframes - black and white, only the main elements, on digital prototype tools - Example here \url{https://miro.medium.com/max/1400/1*eGSfaxwYbxUFomYbyO6GWw.png}
\item High-Fidelity: Wireframes with color, styles, graphical details, and micro-interactions - Example here \url{https://miro.medium.com/max/1400/1*Xn0HSKAvhr4TZzC9lN5udw.gif}
\end{itemize}

\textbf{Deliverables:}

\begin{itemize}
\item Picture of a paper prototype
\item A Figma file with the wireframes (mid-fi)
\item Upload the documents on Github, in a zip folder named "ProjectTitle\_Name\_FirstName", with all project deliverables as follows: "Name\_FirstName\_DeliverableName\_Date\_VersionNumber".  For example, the first deliverable can be named as "Doe\_John\_MidFiPrototype\_05242024\_V1".
\item Add a title within all written documents.
\item Pay attention to the global aspect of the deliverables. They must be clear, simple and easy to read. You can get inspiration on canva.com but don't overload your design with too much details!
\end{itemize}

\textbf{Resources:}

\begin{itemize}
\item Example: Paper Prototyping - The Practical Beginners guide \url{https://www.uxpin.com/studio/blog/paper-prototyping-the-practical-beginners-guide/}
\item Example: Wellness App \url{https://medium.com/@ebtaskin/story-of-my-very-first-ux-ui-design-project-at-ironhack-bootcamp-955addeb760e}
\item Android Mobile Template \url{https://i.pinimg.com/originals/43/8e/de/438ede24264eb9552b635d34142eac29.png}
\item iPhone Template \url{https://www.cssauthor.com/wp-content/uploads/2014/01/iPhone-5S-Template-PSD.jpg}
\item Figma
\end{itemize}

\textbf{Quotes:}

\begin{itemize}
\item 'If you're not embarrassed by the first version of your product, you've launched too late' Reid Hoffman | LinkedIn
\item To prototype your solution, you'll need a temporary change of philosophy: from perfect to just enough, from long-term quality to temporary simulation.
\item The prototype is meant to answer questions, so keep it focused. You just need a real-looking facade to which customers can react.
\end{itemize}

